% ==============================================================================
% CH10_DOCKER_SERVING_TRITON.TEX — CHƯƠNG 10: ĐÓNG GÓI DOCKER, MODEL SERVING VLLM, TRITON & HẠ TẦNG GPU
% ==============================================================================
% Định dạng: Sách chuyên khảo kinh điển (Classic Engineering Monograph)
% Chuẩn mực: Không dùng hộp đóng khung 4 phía, toán học 5 bước, bài tập đa tầng
% ==============================================================================

\chapter{Đóng Gói Docker, Model Serving vLLM, Triton \& Hạ Tầng GPU}
\label{ch:docker-serving-triton}

Một mô hình AI sở hữu năng lực suy luận phi thường trong sổ tay nghiên cứu (Jupyter Notebook) hoàn toàn vô giá trị nếu nó không thể được đóng gói thành một dịch vụ có thể mở rộng (Scalable Service), kiểm soát được tài nguyên phần cứng, và duy trì cam kết chất lượng dịch vụ (SLA) ổn định trước hàng nghìn lượt truy cập đồng thời.

Đưa mô hình AI vào môi trường sản xuất không đơn thuần là bọc một hàm dự đoán bên trong framework FastAPI. Người Kỹ sư AI hiện đại bắt buộc phải làm chủ các tầng trừu tượng sâu nhất của hệ thống: Từ các cơ chế cô lập của nhân Linux (\textbf{Namespaces}, \textbf{cgroups}, \textbf{OverlayFS}), kiến trúc ảo hóa nhân GPU của \textbf{NVIDIA Container Toolkit}, giải phẫu các động cơ phục vụ suy luận chuyên dụng như \textbf{vLLM} và \textbf{Triton Inference Server}, đến chiến lược phân rã ma trận \textbf{Tensor Parallelism (TP)} qua cầu nối siêu tốc \textbf{NVLink} và cơ chế tự động co giãn cụm (\textbf{Autoscaling}) trên Kubernetes dựa trên các số liệu đo lường thời gian thực.

---

\section{Bản Chất Hệ Điều Hành Của Container: Namespaces, Cgroups \& OverlayFS}
\label{sec:linux-container-internals}

Trong tư duy của nhiều người mới bắt đầu, container thường bị hiểu lầm là một ``máy ảo siêu nhẹ''. Đây là một nhận thức sai lầm tai hại. \textbf{Container không phải là một máy ảo}. Bên trong container không có nhân hệ điều hành riêng, không có BIOS ảo, và không có phần cứng ảo hóa. Một container thực chất chỉ là \textbf{một tiến trình Linux thông thường} đang chạy trực tiếp trên nhân của máy chủ chủ (Host Kernel), nhưng được hệ điều hành áp đặt hai cơ chế kiểm soát tối cao: \textbf{Linux Namespaces} (để cô lập tầm nhìn) và \textbf{Control Groups -- cgroups} (để giới hạn tài nguyên).

\subsection{Hệ Thống Linux Namespaces: Ảo Ảnh Không Gian Riêng Biệt}
Linux Namespaces tạo ra ảo ảnh rằng tiến trình đang sở hữu toàn bộ hệ điều hành bằng cách giới hạn những gì tiến trình đó có thể nhìn thấy:

\begin{table}[ht]
\centering
\small
\caption{Các không gian tên (Namespaces) cốt lõi của Linux định hình nên Container.}
\label{tab:linux-namespaces}
\begin{tabular}{lll}
\toprule
\textbf{Namespace} & \textbf{Cờ khởi tạo} & \textbf{Tài nguyên được cô lập} \\
\midrule
\textbf{PID} & \code{CLONE_NEWPID} & Cây tiến trình. Tiến trình chính nhìn thấy mình là PID 1. \\
\textbf{NET} & \code{CLONE_NEWNET} & Giao diện mạng, bảng định tuyến, iptables, cổng socket riêng. \\
\textbf{MNT} & \code{CLONE_NEWNS}  & Cây thư mục hệ thống tệp và các điểm gắn kết (Mount points). \\
\textbf{IPC} & \code{CLONE_NEWIPC} & Hàng đợi tin nhắn POSIX/SysV, bộ nhớ chia sẻ (\code{/dev/shm}). \\
\textbf{UTS} & \code{CLONE_NEWUTS} & Tên máy chủ (Hostname) và miền NIS. \\
\textbf{USER}& \code{CLONE_NEWUSER}& Ánh xạ định danh người dùng (UID/GID) trong và ngoài host. \\
\bottomrule
\end{tabular}
\end{table}

Khi một dịch vụ suy luận AI chạy bên trong container, nó có \code{NET Namespace} riêng biệt với cổng mạng độc lập, và \code{PID Namespace} riêng nơi tiến trình máy chủ web đóng vai trò PID 1. Nếu tiến trình PID 1 bị sập (Crash) hoặc nhận tín hiệu dừng (\code{SIGTERM}), nhân Linux sẽ tự động quét dọn và tiêu diệt toàn bộ các tiến trình con bên trong namespace đó.

\subsection{Control Groups (cgroups v1 vs v2): Cương Tỏa Tài Nguyên Phần Cứng}
Trong khi Namespaces quyết định tiến trình \textbf{nhìn thấy gì}, thì cgroups quyết định tiến trình \textbf{được phép dùng bao nhiêu tài nguyên}.

\paragraph{1. Giới hạn bộ nhớ và cơ chế OOM Killer}:
Bộ điều khiển bộ nhớ (\code{memory controller}) giám sát tổng dung lượng RAM (RSS + Page Cache) của tất cả các tiến trình trong container. Trong cgroups v2, ngưỡng giới hạn tối đa được định nghĩa tại tệp \code{memory.max}. Khi mô hình ngôn ngữ lớn nạp trọng số hoặc phình to bộ nhớ kích hoạt vượt quá ngưỡng này, nhân Linux kích hoạt cơ chế \textbf{Out-Of-Memory (OOM) Killer}: nó tìm kiếm tiến trình có điểm \code{oom_score} cao nhất bên trong cgroup và gửi tín hiệu \code{SIGKILL} (mã thoát 137) để tiêu diệt tiến trình ngay lập tức nhằm bảo vệ an toàn cho máy chủ vật lý.

\paragraph{2. Giới hạn CPU và thảm họa CPU Throttling}:
Bộ điều khiển CPU quản lý thời gian sử dụng nhân vi xử lý thông qua thuật toán lập lịch công bằng \textbf{Completely Fair Scheduler (CFS)}. Hai tham số cốt tử:
\begin{itemize}
  \item \code{cpu.cfs_period_us}: Chu kỳ tính toán quota (thường mặc định là $100.000\text{ }\mu\text{s} = 100\text{ ms}$).
  \item \code{cpu.cfs_quota_us}: Tổng thời gian CPU mà các luồng trong container được phép chạy trong mỗi chu kỳ.
\end{itemize}
Ví dụ: Nếu ta cấu hình Docker \code{--cpus="2.0"}, nhân Linux sẽ thiết lập \code{quota = 200.000} trên \code{period = 100.000}.
\begin{note}
\textbf{Thảm họa CPU Throttling trong dịch vụ AI}: Một máy chủ web bất đồng bộ như FastAPI / Uvicorn xử lý hàng trăm kết nối đồng thời qua Event Loop. Nếu các luồng worker tiêu thụ hết toàn bộ hạn mức $200\text{ ms}$ chỉ trong $30\text{ ms}$ đầu tiên của chu kỳ, nhân Linux sẽ \textbf{cưỡng chế đóng băng (Throttle)} toàn bộ container trong $70\text{ ms}$ còn lại! Hậu quả là độ trễ P99 của API tăng vọt hàng trăm mili-giây một cách bí ẩn, dù mức sử dụng CPU trung bình theo dõi qua lệnh \code{top} có vẻ rất thấp.
\end{note}

\subsection{Hệ Thống Tệp Nhiều Tầng OverlayFS \& Hiệu Ứng Copy-on-Write}
Docker sử dụng trình điều khiển lưu trữ \textbf{OverlayFS} để xếp chồng các lớp ảnh (Image Layers) lên nhau:
\begin{itemize}
  \item \textbf{LowerDir}: Các tầng ảnh chỉ đọc (Read-Only Layers), chứa hệ điều hành cơ sở (Ubuntu), thư viện CUDA và mã nguồn ứng dụng. Các container khác nhau có thể dùng chung các tầng này mà không tốn thêm dung lượng đĩa.
  \item \textbf{UpperDir}: Tầng ghi chép của riêng từng container (Read-Write Container Layer). Mọi tệp mới được tạo ra hoặc chỉnh sửa sẽ lưu tại đây.
  \item \textbf{MergedDir}: Thư mục thống nhất (Union Mount) hiển thị cấu trúc thư mục hoàn chỉnh cho tiến trình bên trong container.
\end{itemize}

\begin{figure}[ht]
\centering
\small
\begin{verbatim}
+--------------------------------------------------------------------------+
| MergedDir: /app (Góc nhìn thống nhất của ứng dụng AI)                     |
+--------------------------------------------------------------------------+
| UpperDir: /var/lib/docker/overlay2/.../diff (Tầng Ghi/Đọc của Container) |
+--------------------------------------------------------------------------+
| LowerDir: Layer 3 - PyTorch & vLLM Packages (Read-Only)                  |
|           Layer 2 - CUDA Runtime & CuDNN (Read-Only)                     |
|           Layer 1 - Ubuntu 22.04 Base Image (Read-Only)                  |
+--------------------------------------------------------------------------+
\end{verbatim}
\caption{Cấu trúc phân tầng lưu trữ của hệ thống tệp OverlayFS.}
\label{fig:overlayfs-layers}
\end{figure}

\paragraph{Cơ chế Sao chép khi Ghi (Copy-on-Write -- CoW)}:
Nếu một tiến trình cố gắng sửa đổi một tệp tồn tại ở \code{LowerDir}, nhân Linux trước hết phải thực hiện thao tác sao chép toàn bộ tệp đó từ \code{LowerDir} lên \code{UpperDir} rồi mới thực hiện phép ghi.
\textit{Cạm bẫy kỹ nghệ}: Nếu kỹ sư vô tình tải tệp trọng số mô hình lớn ($20\text{ GB}$) vào tầng của container thay vì gắn qua thư mục ngoài (\textbf{Mounted Volume}), thao tác CoW sẽ làm tiêu tốn gấp đôi dung lượng đĩa và làm nghẽn toàn bộ bus I/O của máy chủ, khiến container bị sập ổ đĩa (\code{No space left on device})!

---

\section{Ảo Hóa Phần Cứng GPU \& Kiến Trúc NVIDIA Container Toolkit}
\label{sec:gpu-container-virtualization}

\subsection{Tại Sao Docker Mặc Định Hoàn Toàn Bị Mù Trước GPU?}
Nhân Linux tiêu chuẩn quản lý card đồ họa như một thiết bị ngoại vi đặc biệt kết nối qua bus PCIe. Mọi hoạt động giao tiếp phần cứng phải đi qua các tệp ký tự đặc biệt (Device Character Nodes) trong thư mục \code{/dev/}:
\begin{itemize}
  \item \code{/dev/nvidia0}, \code{/dev/nvidia1}, ... : Các chip xử lý GPU vật lý.
  \item \code{/dev/nvidiactl} : Tệp điều khiển tổng thể driver.
  \item \code{/dev/nvidia-uvm} : Mô-đun quản lý bộ nhớ ảo hợp nhất (Unified Virtual Memory).
\end{itemize}
Khi Docker tạo ra một container mới, \code{Device Cgroup} của container mặc định chặn đứng quyền truy cập vào tất cả các tệp thiết bị ký tự ngoại vi này. Do đó, nếu bạn chạy lệnh \code{nvidia-smi} bên trong một container Docker tiêu chuẩn, hệ thống sẽ báo lỗi kinh điển: \code{Failed to initialize NVML: Unknown Error}.

\subsection{Kiến Trúc Phân Tầng Của NVIDIA Container Toolkit}
Để giải quyết bài toán cấp phát phần cứng an toàn, NVIDIA phát triển bộ công cụ \textbf{NVIDIA Container Toolkit} gồm ba thành phần kiến trúc phối hợp:

\begin{figure}[ht]
\centering
\footnotesize
\begin{verbatim}
+--------------------------------------------------------------------------+
| USER-SPACE CONTAINER                                                     |
| PyTorch / vLLM Application                                               |
|   | (Goi ham CUDA Runtime API)                                           |
| cudart64.so (Dóng gói sẵn bên trong Docker Image)                        |
|   | (Goi ham CUDA Driver API)                                            |
| libcuda.so.1 (Driver Stub - Ánh xạ động từ Host vào)                     |
+--------------------------------------------------------------------------+
                                    | (Giao tiếp qua POSIX ioctl syscalls)
+--------------------------------------------------------------------------+
| NVIDIA CONTAINER TOOLKIT (OCI Prestart Hook / CDI)                       |
| 1. Ánh xạ các device nodes: /dev/nvidia* và /dev/nvidia-uvm              |
| 2. Cấu hình Device Cgroup cho phép tiến trình truy cập GPU               |
| 3. Bind-mount các thư viện driver từ Host vào Container (/usr/lib/...)   |
+--------------------------------------------------------------------------+
                                    |
+--------------------------------------------------------------------------+
| HOST KERNEL SPACE                                                        |
| nvidia.ko (Kernel Driver Module) ---> Phần cứng GPU vật lý (A100 / H100) |
+--------------------------------------------------------------------------+
\end{verbatim}
\caption{Kiến trúc phân tầng tương tác phần cứng của NVIDIA Container Toolkit.}
\label{fig:nvidia-toolkit-architecture}
\end{figure}

\begin{enumerate}
  \item \textbf{Cơ chế OCI Prestart Hook}: Khi khởi tạo container với cờ \code{--gpus all}, cấu hình OCI kích hoạt hook tiền khởi động. Công cụ \code{nvidia-container-cli} sẽ cấu hình Device Cgroup và tự động gắn kết các thư viện driver từ máy chủ vào container.
  \item \textbf{Tính di động của Docker Image}: Ảnh Docker của ứng dụng AI \textbf{tuyệt đối không chứa NVIDIA Driver}, mà chỉ chứa \textbf{CUDA Toolkit / CUDA Runtime}. Các thư viện driver tầng thấp (\code{libcuda.so}, \code{libnvidia-ml.so}) được nạp động từ máy chủ khi container khởi chạy, bảo đảm tính tương thích cao.
\end{enumerate}

\paragraph{Phân biệt CUDA Driver Version và CUDA Runtime Version}:
Nhiều kỹ sư thắc mắc khi lệnh \code{nvidia-smi} trên máy chủ báo CUDA Version 12.6, trong khi PyTorch lại in ra giá trị phiên bản \code{12.1}:
\begin{itemize}
  \item \textbf{CUDA Driver}: Do gói cài đặt driver của máy chủ quyết định. Nó quy định phiên bản CUDA tối đa mà phần cứng và driver có thể hỗ trợ.
  \item \textbf{CUDA Runtime}: Là thư viện được biên dịch kèm theo ứng dụng (hoặc gói PyTorch wheel).
  \item \textbf{Nguyên lý tương thích tiến (Forward Compatibility)}: Ứng dụng chạy an toàn và ổn định khi và chỉ khi $\text{CUDA Runtime Version} \le \text{CUDA Driver Version}$.
\end{itemize}

\subsection{Cạm Bẫy Sập Bộ Nhớ Chia Sẻ IPC /dev/shm}
Lỗi phổ biến nhất khiến các container huấn luyện hoặc phục vụ suy luận LLM bị sập với lỗi \code{SIGBUS} (Bus error) là do \textbf{dung lượng bộ nhớ chia sẻ IPC (Shared Memory) mặc định của Docker quá nhỏ}:
\begin{itemize}
  \item Mặc định, Docker chỉ cấp phát đúng \textbf{$64\text{ MB}$} cho phân vùng \code{/dev/shm}.
  \item Khi PyTorch DataLoader sử dụng nhiều luồng (\code{num_workers > 0}) hoặc khi vLLM sử dụng Ray để giao tiếp giữa các GPU trong Tensor Parallelism, các tiến trình chia sẻ tensor qua bộ nhớ chia sẻ này.
  \item Bộ đệm $64\text{ MB}$ bị tràn ngay lập tức trong vài mili-giây đầu tiên, khiến hệ điều hành gửi tín hiệu \code{SIGBUS} giết chết tiến trình ngay lập tức.
  \item \textbf{Biện pháp khắc phục bắt buộc}: Luôn khởi chạy container AI với tham số \code{--shm-size=16g} (hoặc \code{--ipc=host})!
\end{itemize}

---

\section{Động Cơ Phục Vụ Suy Luận: vLLM vs Triton Inference Server}
\label{sec:serving-engines-deepdive}

Việc sử dụng các framework web thông thường (FastAPI, Flask) để bọc mô hình học sâu qua thư viện Hugging Face \code{transformers} là một sai lầm chết người trong sản xuất. Do hạn chế của Python GIL (Global Interpreter Lock), thiếu cơ chế quản lý bộ đệm thông minh, và không hỗ trợ ghép lô động, hệ thống sẽ nhanh chóng rơi vào tình trạng tắc nghẽn hàng đợi và sập OOM khi có tải đồng thời.

Để phục vụ hàng nghìn người dùng với độ trễ thấp nhất, ngành công nghiệp sử dụng hai trường phái động cơ phục vụ suy luận chuyên dụng:

\begin{table}[ht]
\centering
\small
\caption{So sánh chuyên sâu giữa vLLM và Triton Inference Server.}
\label{tab:vllm-vs-triton}
\begin{tabular}{p{3.2cm}p{5.5cm}p{5.5cm}}
\toprule
\textbf{Tiêu chí} & \textbf{vLLM} & \textbf{Triton Inference Server} \\
\midrule
\textbf{Trọng tâm kiến trúc} & Tối ưu hóa chuyên sâu cho Large Language Models (LLMs) & Phục vụ đa mô hình doanh nghiệp (Vision, Audio, NLP, LLM, GNN) \\
\textbf{Quản lý bộ nhớ} & PagedAttention (cấp phát theo trang 16 token, zero-waste) & Dynamic Allocator, chia sẻ bộ nhớ IPC giữa các tiến trình \\
\textbf{Cơ chế ghép lô} & Continuous Batching (Iteration-level Scheduling) & Dynamic Batching (Request-level với timeout) \\
\textbf{Backend hỗ trợ} & PyTorch, FlashAttention, vLLM CUDA Kernels & TensorRT, ONNX Runtime, OpenVINO, PyTorch C++, Python \\
\textbf{Xâu chuỗi Pipeline} & Hạn chế (chủ yếu phục vụ đơn mô hình tạo sinh) & BLS (Business Logic Scripting) và Ensemble Scheduler cực mạnh \\
\textbf{Độ phức tạp triển khai}& Rất thấp (chạy dòng lệnh hoặc image Docker có sẵn) & Cao (đòi hỏi định nghĩa cấu hình \code{config.pbtxt} chi tiết) \\
\bottomrule
\end{tabular}
\end{table}

\subsection{Kiến Trúc Triton Inference Server: Quản Trị Pipeline Doanh Nghiệp}
Triton Inference Server (NVIDIA) tổ chức hệ thống theo cấu trúc \textbf{Kho Mô Hình (Model Repository)}:
\begin{verbatim}
model_repository/
|-- audio_preprocess/
|   |-- config.pbtxt
|   \-- 1/model.py (Python Backend)
|-- whisper_encoder/
|   |-- config.pbtxt
|   \-- 1/model.plan (TensorRT Engine)
\-- llm_orchestrator/
    |-- config.pbtxt (Ensemble Pipeline lien ket cac stage)
    \-- 1/
\end{verbatim}

\paragraph{1. Cơ chế Ghép Lô Động (Dynamic Batching)}:
Trong Triton, dynamic batching được điều khiển bởi hai tham số trọng yếu trong tệp \code{config.pbtxt}:
\begin{itemize}
  \item \code{max_batch_size}: Kích thước lô tối đa mà GPU có thể xử lý trong một lần tính toán.
  \item \code{max_queue_delay_microseconds}: Thời gian tối đa mà Triton cho phép giữ request trong hàng đợi để chờ các request khác đến gom thành một lô lớn hơn.
\end{itemize}
Nếu trong khoảng thời gian timeout này số lượng request đạt \code{max_batch_size}, lô được gửi ngay lập tức vào GPU. Nếu hết thời gian timeout mà chưa đủ lô, Triton vẫn gửi số request hiện có đi xử lý để đảm bảo không vi phạm ràng buộc độ trễ (Latency SLA).

\paragraph{2. Thực Thi Đồng Thời Đa Mô Hình (Concurrent Model Execution)}:
Triton cho phép cấu hình tham số \code{instance_group} để khởi tạo nhiều phiên bản (instances) của cùng một mô hình trên một hoặc nhiều GPU:
\begin{lstlisting}[language=bash]
instance_group [
  {
    count: 2
    kind: KIND_GPU
    gpus: [ 0, 1 ]
  }
]
\end{lstlisting}
Cấu hình trên tạo ra 2 thực thể suy luận trên GPU 0 và 2 thực thể trên GPU 1, cho phép xử lý đồng thời 4 luồng dự đoán độc lập mà không bị chặn bởi tiến trình khác!

\paragraph{3. Business Logic Scripting (BLS) \& Pipeline Ensemble}:
Trong các hệ thống AI thực tế, dữ liệu thô (như văn bản chưa tách từ, file âm thanh WAV hoặc ảnh JPEG) cần trải qua nhiều bước tiền xử lý trước khi đưa vào mô hình chính. Thay vì gửi dữ liệu qua lại giữa các microservices qua mạng HTTP làm phát sinh độ trễ khổng lồ, Triton cho phép xâu chuỗi các mô hình trực tiếp trong bộ nhớ chia sẻ máy chủ thông qua \textbf{BLS (Business Logic Scripting)}:
\begin{equation}
\text{Audio RAW} \xrightarrow[\text{Zero-Copy Memory}]{\text{Python Preprocess}} \text{Spectrogram} \xrightarrow[\text{TensorRT GPU}]{\text{Whisper}} \text{Tokens} \xrightarrow[\text{vLLM Engine}]{\text{Llama-3}} \text{Response}
\end{equation}
Toàn bộ chuỗi xử lý diễn ra trên cùng một không gian bộ nhớ máy chủ với độ trễ chuyển giao tiệm cận $0\text{ ms}$!

---

\section{Song Song Hóa Mô Hình: Tensor Parallelism vs Pipeline Parallelism}
\label{sec:model-parallelism}

Khi kích thước mô hình vượt quá dung lượng VRAM của một GPU vật lý đơn lẻ (ví dụ: mô hình Llama-3-70B ở định dạng FP16 chiếm $140\text{ GB}$, trong khi GPU Nvidia A100 chỉ có $80\text{ GB}$ VRAM), chúng ta bắt buộc phải phân chia mô hình ra nhiều GPU cùng chạy. Có hai chiến lược phân chia kiến trúc chính:

\subsection{Tensor Parallelism (TP) -- Phân Rã Ma Trận Trọng Số}
Được phát minh trong công trình Megatron-LM (Shoeybi et al., NVIDIA, 2019), Tensor Parallelism chia nhỏ các phép tính nhân ma trận bên trong từng tầng Transformer ra nhiều GPU cùng tính toán song song:

\paragraph{1. Tầng Chiếu Đầu Vào / MLP Cột (Column-Parallel Linear)}:
Giả sử tầng nơ-ron thực hiện phép nhân giữa véc-tơ kích hoạt $\mathbf{x} \in \mathbb{R}^{1 \times d_{\text{in}}}$ và ma trận trọng số $\mathbf{W} \in \mathbb{R}^{d_{\text{in}} \times d_{\text{out}}}$.
Nếu hệ thống có $N$ GPU (ví dụ $N=2$), ta chia ma trận $\mathbf{W}$ theo chiều dọc (cột) thành $N$ phần bằng nhau:
\begin{equation}
\mathbf{W} = \begin{bmatrix} \mathbf{W}_1 & \mathbf{W}_2 \end{bmatrix} \quad \text{với } \mathbf{W}_i \in \mathbb{R}^{d_{\text{in}} \times \frac{d_{\text{out}}}{2}}
\end{equation}
Mỗi GPU giữ véc-tơ đầu vào $\mathbf{x}$ giống nhau và độc lập tính toán một nửa số cột của đầu ra:
\begin{equation}
\mathbf{y}_1 = \mathbf{x} \mathbf{W}_1, \qquad \mathbf{y}_2 = \mathbf{x} \mathbf{W}_2
\end{equation}
Giai đoạn này \textbf{hoàn toàn không cần bất kỳ sự giao tiếp nào giữa các GPU}!

\paragraph{2. Tầng Chiếu Đầu Ra / MLP Hàng (Row-Parallel Linear)}:
Ở tầng tiếp theo, để đưa véc-tơ về kích thước ban đầu, ma trận trọng số thứ hai $\mathbf{V} \in \mathbb{R}^{d_{\text{out}} \times d_{\text{final}}}$ được chia theo chiều ngang (hàng):
\begin{equation}
\mathbf{V} = \begin{bmatrix} \mathbf{V}_1 \\ \mathbf{V}_2 \end{bmatrix} \quad \text{với } \mathbf{V}_i \in \mathbb{R}^{\frac{d_{\text{out}}}{2} \times d_{\text{final}}}
\end{equation}
Mỗi GPU nhân kết quả cục bộ của mình với nửa ma trận hàng tương ứng:
\begin{equation}
\mathbf{z}_1 = \mathbf{y}_1 \mathbf{V}_1, \qquad \mathbf{z}_2 = \mathbf{y}_2 \mathbf{V}_2
\end{equation}
Theo quy tắc đại số tuyến tính:
\begin{equation}
\mathbf{z} = \mathbf{y} \mathbf{V} = \begin{bmatrix} \mathbf{y}_1 & \mathbf{y}_2 \end{bmatrix} \begin{bmatrix} \mathbf{V}_1 \\ \mathbf{V}_2 \end{bmatrix} = \mathbf{y}_1 \mathbf{V}_1 + \mathbf{y}_2 \mathbf{V}_2 = \mathbf{z}_1 + \mathbf{z}_2
\end{equation}
Đến lúc này, hai GPU bắt buộc phải thực hiện thao tác đồng bộ cộng tổng: \textbf{All-Reduce Sum}!

\begin{tradeoffbox}[Giao Tiếp All-Reduce \& Ràng Buộc Phần Cứng Khắt Khe Của Tensor Parallelism]
Tại mỗi tầng Transformer, Tensor Parallelism yêu cầu đúng \textbf{hai lần thao tác All-Reduce} (một lần sau khối Multi-Head Attention, và một lần sau khối MLP).
Một mô hình 80 tầng như Llama-3-70B thực hiện:
\begin{equation*}
80 \times 2 = 160\text{ lần All-Reduce cho MỖI MỘT TOKEN được sinh ra!}
\end{equation*}
Nếu độ trễ mỗi lần giao tiếp là $10\text{ micro-giây}$, tổng thời gian mạng cho 1 token là $1.6\text{ ms}$.
Nếu chạy qua cáp mạng LAN thông thường hoặc bus PCIe chậm chạp, thời gian giao tiếp mạng sẽ nuốt chửng toàn bộ thời gian tính toán!
Do đó:
\begin{itemize}
  \item \textbf{Quy tắc vàng}: Tensor Parallelism \textbf{bắt buộc chỉ được chạy trong cùng một máy chủ vật lý (Intra-Node)} nơi các GPU được liên kết trực tiếp bằng cầu nối siêu tốc \textbf{NVIDIA NVLink} (băng thông từ $600\text{ GB/s}$ đến $3.200\text{ GB/s}$).
  \item Tuyệt đối không bao giờ cấu hình Tensor Parallelism xuyên qua các máy chủ khác nhau qua mạng Ethernet thông thường!
\end{itemize}
\end{tradeoffbox}

\subsection{Pipeline Parallelism (PP) -- Phân Rã Theo Chiều Sâu Các Tầng}
Nếu một cụm máy chủ gồm nhiều node mạng và mô hình quá khổng lồ (như Llama-3-405B), chúng ta sử dụng \textbf{Pipeline Parallelism (PP)}:
\begin{itemize}
  \item Chia 80 tầng Transformer thành các khối: Node 1 giữ tầng 1--20, Node 2 giữ tầng 21--40, Node 3 giữ tầng 41--60, Node 4 giữ tầng 61--80.
  \item Giao tiếp giữa các Node chỉ xảy ra ở ranh giới giữa các khối (Node 1 tính xong tầng 20 thì gửi Activation qua mạng cho Node 2).
  \item Nhược điểm lớn: Hiện tượng \textbf{Bong bóng Pipeline (Pipeline Bubbles)} khi các node phía sau phải ngồi chờ dữ liệu từ các node phía trước. Các thuật toán lập lịch hiện đại như \textbf{1F1B (One-Forward-One-Backward)} được áp dụng để xen kẽ các micro-batch nhằm giảm thiểu thời gian chờ rảnh rỗi này.
\end{itemize}

---

\section{Vận Hành Sản Xuất: Dockerfile Đa Tầng \& Tự Động Co Giãn KEDA Kubernetes}
\label{sec:production-docker-keda}

\subsection{Kỹ Nghệ Tối Ưu Hóa Dockerfile Đa Tầng Cho Dịch Vụ AI}
Một image Docker phục vụ sản xuất phải thỏa mãn ba tiêu chí khắt khe: Kích thước nhỏ nhất có thể, thời gian khởi động nhanh nhất, và bảo mật tuyệt đối không để lộ đặc quyền root.

\begin{lstlisting}[language=bash, caption={Dockerfile đa tầng tối ưu hóa cho dịch vụ phục vụ vLLM production.}]
# ==============================================================================
# STAGE 1: Builder - Bien dich thu vien C++/CUDA va cai dat dependencies
# ==============================================================================
FROM nvidia/cuda:12.4.1-devel-ubuntu22.04 AS builder

ENV DEBIAN_FRONTEND=noninteractive \
    PIP_NO_CACHE_DIR=1 \
    PYTHONUNBUFFERED=1

RUN apt-get update && apt-get install -y --no-install-recommends \
    python3.11 \
    python3.11-dev \
    python3.11-venv \
    python3-pip \
    build-essential \
    git \
    && rm -rf /var/lib/apt/lists/*

# Khoi tao moi truong ao cach ly
RUN python3.11 -m venv /opt/venv
ENV PATH="/opt/venv/bin:$PATH"

# Sao chep danh sach goi phu thuoc truoc de tan dung Docker Cache Layer
COPY requirements.txt /tmp/requirements.txt
RUN pip install --upgrade pip setuptools wheel && \
    pip install -r /tmp/requirements.txt

# ==============================================================================
# STAGE 2: Runtime Base - Chi chua CUDA Runtime va thu vien can thiet
# ==============================================================================
FROM nvidia/cuda:12.4.1-runtime-ubuntu22.04 AS runtime

ENV DEBIAN_FRONTEND=noninteractive \
    PYTHONUNBUFFERED=1 \
    PATH="/opt/venv/bin:$PATH" \
    HF_HOME=/app/cache/huggingface

WORKDIR /app

RUN apt-get update && apt-get install -y --no-install-recommends \
    python3.11 \
    curl \
    ca-certificates \
    && rm -rf /var/lib/apt/lists/*

# Sao chep toan bo thu vien da bien dich tu Builder Stage (Loai bo sach se trinh bien dich)
COPY --from=builder /opt/venv /opt/venv

# Tao tai khoan non-root voi UID co dinh de bao mat Kubernetes
RUN useradd -m -u 10001 aiuser && \
    mkdir -p /app/cache && \
    chown -R aiuser:aiuser /app

USER aiuser
EXPOSE 8000

HEALTHCHECK --interval=30s --timeout=5s --start-period=60s --retries=3 \
  CMD curl -f http://localhost:8000/health || exit 1

ENTRYPOINT ["python3", "-m", "vllm.entrypoints.openai.api_server"]
CMD ["--model", "meta-llama/Meta-Llama-3-8B-Instruct", "--port", "8000", "--gpu-memory-utilization", "0.90"]
\end{lstlisting}

\subsection{Autoscaling Cho Tác Vụ AI: Tại Sao HPA Truyền Thống Thất Bại?}
Trên các cụm Kubernetes thông thường, bộ tự động co giãn ngang (\textbf{Horizontal Pod Autoscaler -- HPA}) điều khiển số lượng bản sao (Replicas) dựa trên mức sử dụng CPU hoặc RAM:
\begin{equation}
\text{CPU Utilization} = \frac{\text{Tổng CPU tiêu thụ}}{\text{CPU Request}} > 80\% \implies \text{Scale Up Pod}
\end{equation}
Tuy nhiên, đối với dịch vụ phục vụ LLM, HPA truyền thống hoàn toàn bị vô hiệu hóa vì hai lý do cốt tử:
\begin{enumerate}
  \item \textbf{VRAM luôn ở mức 90\%--95\% cố định}: Ngay khi vừa khởi động, vLLM đã chiếm dụng trước (Pre-allocate) toàn bộ dung lượng VRAM còn trống cho PagedAttention KV Cache. Do đó, chỉ số sử dụng bộ nhớ luôn chạm ngưỡng $95\%$ ngay cả khi không có người dùng nào gửi request!
  \item \textbf{CPU thấp trong khi hàng đợi tắc nghẽn}: Khi lưu lượng request tăng đột biến, GPU bị quá tải và hàng trăm request phải xếp hàng chờ trong hàng đợi (\code{waiting queue}). Tuy nhiên, tiến trình CPU chỉ làm nhiệm vụ điều phối nhẹ nên mức sử dụng CPU chỉ dao động $20\% - 30\%$. HPA hiểu lầm hệ thống vẫn đang rảnh rỗi và từ chối mở rộng pod, dẫn tới thảm họa sập dịch vụ!
\end{enumerate}

\paragraph{Giải pháp cứu cánh: KEDA (Kubernetes Event-driven Autoscaling)}:
Chúng ta sử dụng KEDA để co giãn pod dựa trên các số liệu đo lường thời gian thực từ endpoint Prometheus của vLLM:
\begin{itemize}
  \item \textbf{Chỉ số số lượng request đang chờ (\code{vllm:num_requests_waiting})}: Số lượng yêu cầu bị tắc nghẽn trong hàng đợi do thiếu slot KV Cache.
  \item \textbf{Chỉ số chiếm dụng KV Cache (\code{vllm:gpu_cache_usage_factor})}: Tỷ lệ phần trăm các block bộ nhớ PagedAttention đang được sử dụng.
\end{itemize}

Công thức xác định số lượng pod mục tiêu theo tải hàng đợi:
\begin{equation}
R_{\text{target}} = \max\left( R_{\min}, \min\left( R_{\max}, \left\lceil \frac{Q_{\text{waiting}} + B_{\text{active}}}{C_{\text{pod}}} \right\rceil \right) \right)
\end{equation}
Trong đó $Q_{\text{waiting}}$ là tổng số request đang chờ, $B_{\text{active}}$ là số request đang sinh token, và $C_{\text{pod}}$ là dung lượng xử lý đồng thời tối ưu của mỗi pod (Concurrency capacity).

---

\section{Ví Dụ Giải Mẫu Tính Toán Số Học Bằng Tay}
\label{sec:serving-worked-examples}

\begin{example}[Bài Toán Băng Thông Giao Tiếp All-Reduce Trong Tensor Parallelism]
Xét cụm phục vụ mô hình \textbf{Llama-3-70B} sử dụng Tensor Parallelism với $TP = 4$ GPU:
\begin{itemize}
  \item Số chiều ẩn: $d_{\text{model}} = 8.192$.
  \item Số tầng Transformer: $n_{\text{layers}} = 80$.
  \item Kích thước lô phục vụ pha Decode: $B = 32$ người dùng đồng thời.
  \item Định dạng số thực nửa độ chính xác: FP16 ($b = 2$ bytes).
  \item Mỗi tầng Transformer có đúng $2$ phép tính All-Reduce (sau Attention và sau MLP).
\end{itemize}

Hãy tính toán:
\begin{enumerate}
  \item Dung lượng của véc-tơ kích hoạt cần truyền trong mỗi phép tính All-Reduce.
  \item Tổng lượng dữ liệu mà mỗi GPU phải truyền nhận qua mạng trong thuật toán Ring All-Reduce cho đúng $1$ token được sinh ra.
  \item Thời gian trễ mạng thuần túy cho 1 token nếu sử dụng cầu nối \textbf{NVIDIA NVLink} (Băng thông: $900\text{ GB/s}$).
  \item Thời gian trễ mạng thuần túy cho 1 token nếu sử dụng khe cắm \textbf{PCIe 4.0 x16} truyền thống (Băng thông thực tế: $32\text{ GB/s}$). So sánh và bình luận kết quả.
\end{enumerate}
\end{example}

\paragraph{Lời giải chi tiết từng bước}:
\begin{itemize}
  \item[1)] Kích thước tensor truyền thông trong mỗi phép All-Reduce:
  \begin{equation*}
  S = B \times d_{\text{model}} \times b = 32 \times 8.192 \times 2\text{ Bytes} = 524.288\text{ Bytes} = 512\text{ KiB} = 0.5\text{ MiB}
  \end{equation*}

  \item[2)] Trong thuật toán Ring All-Reduce với $N = 4$ GPU, mỗi GPU truyền và nhận đúng:
  \begin{equation*}
  D_{\text{All-Reduce}} = 2 \times \left( \frac{N - 1}{N} \right) \times S = 2 \times \frac{3}{4} \times 512\text{ KiB} = 768\text{ KiB}
  \end{equation*}
  Tổng số lần All-Reduce để sinh ra đúng $1$ token qua 80 tầng:
  \begin{equation*}
  N_{\text{ops}} = 80 \times 2 = 160\text{ lần}
  \end{equation*}
  Tổng lượng dữ liệu mỗi GPU trao đổi cho 1 token:
  \begin{equation*}
  D_{\text{total}} = 160 \times 768\text{ KiB} = 122.880\text{ KiB} = 120\text{ MiB} \approx 0.1258\text{ GB}
  \end{equation*}

  \item[3)] Thời gian trễ mạng qua NVIDIA NVLink ($900\text{ GB/s}$):
  \begin{equation*}
  t_{\text{NVLink}} = \frac{0.1258\text{ GB}}{900\text{ GB/s}} \approx 0.0001398\text{ giây} \approx 0.14\text{ ms}
  \end{equation*}
  Độ trễ $0.14\text{ ms}$ là không đáng kể so với thời gian tính toán sinh token ($15 - 25\text{ ms}$).

  \item[4)] Thời gian trễ mạng qua PCIe 4.0 x16 ($32\text{ GB/s}$):
  \begin{equation*}
  t_{\text{PCIe}} = \frac{0.1258\text{ GB}}{32\text{ GB/s}} \approx 0.003931\text{ giây} \approx 3.93\text{ ms}
  \end{equation*}
  Thời gian mạng tăng vọt gần \textbf{$28$ lần}! Giao tiếp ngốn gần $4\text{ ms}$ trên mỗi token, làm giảm thông lượng toàn hệ thống nghiêm trọng.
\end{itemize}

---

\begin{example}[Định Lượng Dynamic Batching \& Tối Ưu Hóa Độ Trễ SLA]
Một dịch vụ phân loại văn bản AI sử dụng Triton Inference Server với thời gian tính toán trên GPU phụ thuộc tuyến tính vào kích thước lô $b$:
\begin{equation*}
T_{\text{compute}}(b) = 4.0 + 0.5 \cdot b \quad (\text{mili-giây})
\end{equation*}
Lưu lượng yêu cầu gửi đến tuân theo phân phối Poisson với tốc độ trung bình $\lambda = 40\text{ requests/giây}$ (tương đương $0.04\text{ req/ms}$).
Hệ thống cấu hình \code{max_batch_size = 16}.
Cam kết chất lượng dịch vụ (SLA) yêu cầu tổng thời gian xử lý một request không được vượt quá $15.0\text{ ms}$.
Hãy xác định:
\begin{enumerate}
  \item Thời gian tính toán khi chạy đơn lẻ ($b=1$) và khi chạy đầy lô ($b=16$).
  \item Ngân sách thời gian tối đa cho phép giữ request trong hàng đợi (\code{max_queue_delay_microseconds}) để không vi phạm SLA $15.0\text{ ms}$ ngay cả khi lô đạt kích thước cực đại 16.
  \item Với thời gian chờ tối đa đó, số lượng request trung bình gom được trong một chu kỳ là bao nhiêu?
\end{enumerate}
\end{example}

\paragraph{Lời giải chi tiết}:
\begin{itemize}
  \item[1)] Tính thời gian thực thi:
  \begin{itemize}
    \item Khi $b = 1$: $T_{\text{compute}}(1) = 4.0 + 0.5 \times 1 = 4.5\text{ ms}$.
    \item Khi $b = 16$: $T_{\text{compute}}(16) = 4.0 + 0.5 \times 16 = 4.0 + 8.0 = 12.0\text{ ms}$.
  \end{itemize}

  \item[2)] Tổng thời gian một request phải trải qua là:
  \begin{equation*}
  T_{\text{total}} = T_{\text{wait}} + T_{\text{compute}}(b) \le 15.0\text{ ms}
  \end{equation*}
  Trong trường hợp xấu nhất, lô đạt kích thước $b=16$ và tốn $12.0\text{ ms}$ để tính toán. Do đó, thời gian chờ tối đa trong hàng đợi chỉ được phép là:
  \begin{equation*}
  T_{\text{wait\_max}} = 15.0 - 12.0 = 3.0\text{ ms} = 3000\text{ }\mu\text{s}
  \end{equation*}
  Vậy giá trị tối ưu để cấu hình là \code{max_queue_delay_microseconds: 3000}.

  \item[3)] Với $\lambda = 0.04\text{ req/ms}$ và khoảng thời gian chờ $T_{\text{wait}} = 3.0\text{ ms}$, số lượng request trung bình đến trong khoảng thời gian này là:
  \begin{equation*}
  \bar{N} = \lambda \times T_{\text{wait}} = 0.04 \times 3.0 = 0.12\text{ request}
  \end{equation*}
  \textit{Nhận xét}: Với lưu lượng $\lambda = 40\text{ req/s}$, hệ thống rất ít khi gom đủ lô 16 trong vòng $3\text{ ms}$. Để tận dụng tối đa GPU, cần có cơ chế cân bằng tải hoặc tăng lưu lượng qua việc phục vụ đa khách hàng.
\end{itemize}

---

\section{Hiện Thực Hóa: Mã Nguồn Python Chuẩn Kỹ Nghệ}
\label{sec:python-serving-modules}

Dưới đây là các module Python thuần chuẩn mực, tự chứa và không phụ thuộc vào thư viện ngoài, hiện thực hóa bộ mô phỏng phân rã Tensor Parallelism, bộ lập lịch Dynamic Batching, và công cụ tính toán co giãn pod tự động:

\begin{lstlisting}[language=Python, caption={Mô phỏng phép nhân ma trận Column-Parallel trong Tensor Parallelism.}]
import numpy as np

def simulate_tensor_parallelism_linear(
    x: np.ndarray,
    w: np.ndarray,
    tp_size: int = 2
) -> np.ndarray:
    """
    Mo phong phep chieu Column-Parallel Linear trong Tensor Parallelism.
    Chia ma tran W theo chieu cot ra tp_size GPU.
    Kiem chung: Ket qua gop tu cac GPU phai trung khop voi x @ w goc.
    """
    d_in, d_mid = w.shape
    assert d_mid % tp_size == 0, "So chieu an phai chia het cho TP size"
    split_size = d_mid // tp_size

    # 1. Column-Parallel: Chia W thanh cac khoi cot
    w_splits = [w[:, i * split_size:(i + 1) * split_size] for i in range(tp_size)]

    # 2. Moi GPU thuc hien phep nhan cuc bo voi vector x
    local_outputs = [np.matmul(x, w_i) for w_i in w_splits]

    # 3. Gop ket qua trung gian (Concatenate doc theo chieu cot)
    full_intermediate = np.concatenate(local_outputs, axis=-1)
    
    # Kiem tra tinh tuong duong toan hoc
    expected = np.matmul(x, w)
    assert np.allclose(full_intermediate, expected, atol=1e-5), "Loi sai lech so hoc TP!"

    return full_intermediate
\end{lstlisting}

\begin{lstlisting}[language=Python, caption={Bộ lập lịch Dynamic Batch Scheduler hỗ trợ timeout và giới hạn lô.}]
import time
from typing import List, Any, Tuple

class DynamicBatchScheduler:
    """
    Bo lap lich gop lo dong (Dynamic Batch Scheduler):
    Gom cac request trong hang doi thanh mot batch dua tren max_batch_size
    hoac khi cham nguong thoi gian cho max_delay_ms.
    """
    def __init__(self, max_batch_size: int = 8, max_delay_ms: float = 10.0):
        self.max_batch_size = max_batch_size
        self.max_delay_ms = max_delay_ms
        self.queue: List[Tuple[Any, float]] = []  # (item, arrival_time)

    def enqueue(self, item: Any, arrival_time: float = None):
        """Dua mot request moi vao hang doi."""
        if arrival_time is None:
            arrival_time = time.time()
        self.queue.append((item, arrival_time))

    def should_dispatch(self, current_time: float = None) -> bool:
        """Kiem tra xem da den luc ban giao batch di xu ly chua."""
        if not self.queue:
            return False
        if len(self.queue) >= self.max_batch_size:
            return True
        if current_time is None:
            current_time = time.time()
        oldest_arrival = self.queue[0][1]
        elapsed_ms = (current_time - oldest_arrival) * 1000.0
        return elapsed_ms >= self.max_delay_ms

    def dispatch_batch(self) -> List[Any]:
        """Lay ra mot batch request de xu ly tren GPU."""
        batch_count = min(len(self.queue), self.max_batch_size)
        dispatched_items = [self.queue[i][0] for i in range(batch_count)]
        self.queue = self.queue[batch_count:]
        return dispatched_items
\end{lstlisting}

\begin{lstlisting}[language=Python, caption={Động cơ tính toán co giãn Pod tự động KEDA cho AI Serving.}]
import math

class KEDAAutoscalingCalculator:
    """
    Tinh toan so luong Pod can co gian dua tren metrics thoi gian thuc cua vLLM:
    - waiting_requests: So luong request bi tac nghen trong queue
    - active_requests: So luong request dang duoc xu ly
    - pod_capacity: So request dong thoi toi uu ma 1 GPU Pod xu ly duoc
    """
    def __init__(self, min_replicas: int = 1, max_replicas: int = 10, pod_capacity: int = 32):
        self.min_replicas = min_replicas
        self.max_replicas = max_replicas
        self.pod_capacity = pod_capacity

    def compute_replicas(self, waiting_requests: int, active_requests: int) -> int:
        total_load = waiting_requests + active_requests
        if total_load == 0:
            return self.min_replicas

        desired_replicas = math.ceil(total_load / self.pod_capacity)
        clamped_replicas = max(self.min_replicas, min(self.max_replicas, desired_replicas))
        return clamped_replicas
\end{lstlisting}

---

\section{Hạ Tầng Mạng Siêu Tốc Cho Cụm AI: RDMA, RoCEv2, InfiniBand \& Thuật Toán NCCL}
\label{sec:ai-high-performance-networking}

Khi huấn luyện hoặc phục vụ các mô hình ngôn ngữ lớn quy mô hàng chục tỷ tham số, việc tính toán không thể gói gọn trong một máy chủ GPU đơn lẻ. Các kỹ sư bắt buộc phải phân bổ mô hình qua hàng chục, thậm chí hàng ngàn bộ vi xử lý đồ họa thông qua chiến lược song song hóa (\textit{Tensor Parallelism, Pipeline Parallelism, Data Parallelism}).
Lúc này, một thách thức vật lý mang tính sống còn xuất hiện: \textbf{Mạng truyền thông liên máy chủ (Inter-Node Interconnect) trở thành điểm nghẽn cổ chai khốc liệt nhất của toàn bộ hệ thống}.

\subsection{Khởi Điểm Từ Con Số 0: Tại Sao Mạng TCP/IP Truyền Thống Là Điểm Nghẽn Của GPU?}
Để thấu hiểu tại sao các kỹ sư hạ tầng AI không bao giờ sử dụng giao thức mạng TCP/IP thông thường cho giao tiếp giữa các GPU, hãy xem xét hành trình của một gói tin qua ngăn xếp giao vận truyền thống:

\paragraph{Hành trình 5 bước đầy tắc nghẽn của TCP/IP}:
\begin{enumerate}
  \item \textbf{Sao chép từ bộ nhớ GPU sang RAM chủ}: Dữ liệu tensor từ bộ nhớ VRAM của GPU phải được truyền qua bus PCIe để nạp vào bộ nhớ RAM của CPU (Host Memory).
  \item \textbf{Chuyển đổi ngữ cảnh nhân hệ điều hành (Kernel Context Switch)}: Tiến trình người dùng gọi hàm hệ thống (\textit{Syscall} như \code{send()} hoặc \code{write()}), buộc CPU phải chuyển từ User Space sang Kernel Space.
  \item \textbf{Cấp phát bộ nhớ đệm Socket Buffer}: Nhân Linux cấp phát cấu trúc bộ nhớ \code{sk_buff}, sao chép dữ liệu từ bộ đệm người dùng vào bộ đệm của nhân, tính toán mã kiểm tra chẵn lẻ (Checksum), và đóng gói các tiêu đề TCP/IP.
  \item \textbf{Đẩy dữ liệu sang Card Mạng (NIC DMA)}: Trình điều khiển card mạng sử dụng cơ chế Direct Memory Access (DMA) để kéo dữ liệu từ RAM chủ sang bộ nhớ của card mạng và bắn ra cáp quang.
  \item \textbf{Xử lý ngắt phần cứng tại đích}: Máy chủ nhận phải phát sinh hàng loạt ngắt phần cứng (\textit{Hardware Interrupts}) để đánh thức CPU xử lý gói tin và lặp lại toàn bộ quy trình sao chép ngược lên GPU đích.
\end{enumerate}

\paragraph{Nghịch lý độ trễ microsecond vs nanosecond}:
Toàn bộ hành trình trên tạo ra độ trễ mạng một chiều tối thiểu từ $15\mu\text{s}$ đến $40\mu\text{s}$ (microsecond) cùng với việc chiếm dụng $100\%$ công suất của nhiều nhân CPU chỉ để sao chép dữ liệu. Trong khi đó, các lõi Tensor Cores trên card NVIDIA H100 thực hiện một phép nhân ma trận FP8 khổng lồ chỉ trong vòng vài trăm nanoseconds ($0.3\mu\text{s}$). Nếu sử dụng TCP/IP, các GPU đắt giá sẽ phải ngồi rảnh rỗi chờ đợi dữ liệu mạng tới hơn $98\%$ thời gian!

\subsection{Kiến Trúc RDMA (Remote Direct Memory Access) \& Mạng Lossless Ethernet (RoCEv2)}
Để giải phóng hoàn toàn CPU và triệt tiêu độ trễ ngăn xếp mạng, công nghệ đột phá tối cao của hạ tầng AI hiện đại là \term{RDMA (Remote Direct Memory Access)} kết hợp \term{GPUDirect RDMA}.

\begin{figure}[htbp]
\centering
\begin{tikzpicture}[node distance=1.2cm, auto,
    host/.style={rectangle, draw=primarycolor, thick, fill=boxbg, text centered, rounded corners=3pt, minimum height=3.8cm, minimum width=6.0cm},
    comp/.style={rectangle, draw=secondarycolor, thick, fill=accentcolor!15, text centered, rounded corners=2pt, minimum height=0.7cm, minimum width=2.4cm, font=\scriptsize\bfseries}
]
    % Node A
    \node[host] (hostA) {};
    \node[above=0.1cm of hostA.north, font=\small\bfseries\color{primarycolor}] {Máy Chủ GPU A (Sender)};
    \node[comp, at={(hostA.north west)}, xshift=1.5cm, yshift=-0.8cm] (gpuA) {GPU VRAM A};
    \node[comp, at={(hostA.west)}, xshift=1.5cm] (pcieA) {PCIe Switch};
    \node[comp, at={(hostA.south west)}, xshift=1.5cm, yshift=0.8cm] (nicA) {SmartNIC / HCA};
    
    % Node B
    \node[host, right=3.5cm of hostA] (hostB) {};
    \node[above=0.1cm of hostB.north, font=\small\bfseries\color{primarycolor}] {Máy Chủ GPU B (Receiver)};
    \node[comp, at={(hostB.north east)}, xshift=-1.5cm, yshift=-0.8cm] (gpuB) {GPU VRAM B};
    \node[comp, at={(hostB.east)}, xshift=-1.5cm] (pcieB) {PCIe Switch};
    \node[comp, at={(hostB.south east)}, xshift=-1.5cm, yshift=0.8cm] (nicB) {SmartNIC / HCA};

    % Links Node A
    \draw[draw=primarycolor, ultra thick, Stealth-Stealth] (gpuA) -- (pcieA);
    \draw[draw=primarycolor, ultra thick, Stealth-Stealth] (pcieA) -- (nicA);
    % Links Node B
    \draw[draw=primarycolor, ultra thick, Stealth-Stealth] (gpuB) -- (pcieB);
    \draw[draw=primarycolor, ultra thick, Stealth-Stealth] (pcieB) -- (nicB);

    % Direct RDMA network link
    \draw[draw=dangercolor, line width=2.5pt, Stealth-Stealth] (nicA) -- node[below, font=\scriptsize\bfseries\color{dangercolor}] {Cáp Quang 400/800 Gbps (InfiniBand / RoCEv2)} (nicB);
    
    % Bypass indicator
    \node[draw=dangercolor, dashed, thick, fill=white, rounded corners=2pt, font=\tiny\bfseries\color{dangercolor}, above=0.6cm of nicA, xshift=4.2cm] {Zero-Copy \& Kernel Bypass: CPU và RAM hoàn toàn không tham gia!};
\end{tikzpicture}
\caption{Kiến trúc GPUDirect RDMA: Truyền dữ liệu thẳng từ VRAM máy chủ này sang VRAM máy chủ khác}
\label{fig:gpudirect-rdma}
\end{figure}

\paragraph{Hai nguyên lý nền tảng của RDMA}:
\begin{enumerate}
  \item \textbf{Kernel Bypass (Bỏ qua nhân hệ điều hành)}: Thư viện ứng dụng người dùng (như NCCL hoặc UCX) tương tác trực tiếp với bộ điều khiển card mạng (Hardware Queue Pairs) mà không hề gọi System Call.
  \item \textbf{Zero-Copy qua GPUDirect RDMA}: Nhờ sự hỗ trợ của chuẩn PCIe P2P (Peer-to-Peer), card mạng thông minh (SmartNIC như Mellanox ConnectX-7) đọc trực tiếp dữ liệu từ VRAM của GPU qua bus PCIe nội bộ và bắn trực tiếp sang VRAM của GPU từ xa qua mạng cáp quang. CPU chủ và bộ nhớ RAM hoàn toàn được giải phóng.
  Độ trễ truyền mạng giảm từ $30\mu\text{s}$ xuống dưới \textbf{$1.2\mu\text{s}$}, với băng thông chạm ngưỡng phần cứng ($400\text{ Gbps}$ đến $800\text{ Gbps}$).
\end{enumerate}

\paragraph{Mạng không mất gói (Lossless Ethernet qua RoCEv2)}:
Trong khi mạng InfiniBand sử dụng cơ chế điều khiển luồng dựa trên tín dụng cấp phần cứng (\textit{Credit-based Flow Control}) vốn dĩ không bao giờ làm rơi gói tin, thì giao thức mạng phổ thông \term{RoCEv2 (RDMA over Converged Ethernet)} phải chạy trên hạ tầng mạng cáp Ethernet truyền thống.
Giao thức vận chuyển của RDMA không sử dụng cơ chế truyền lại thông minh của TCP (khi mất một gói tin, nó phải truyền lại toàn bộ từ đầu qua cơ chế Go-Back-N, làm suy sụp băng thông tới $99\%$). Do đó, mạng Ethernet dành cho cụm AI bắt buộc phải được cấu hình thành \term{Lossless Ethernet} thông qua hai cơ chế phần cứng:
\begin{itemize}
  \item \textbf{PFC (Priority Flow Control -- IEEE 802.1Qbb)}: Chia băng thông cáp vật lý thành 8 hàng đợi ưu tiên độc lập (Traffic Classes). Lưu lượng RDMA của AI được gán riêng vào hàng đợi có mức ưu tiên cao (thường là Hàng đợi số 3). Khi bộ đệm hàng đợi này trên Switch đạt tới ngưỡng cảnh báo (\textit{PFC Threshold}), Switch lập tức gửi khung tín hiệu \code{PAUSE Frame} ngược về card mạng phát để tạm dừng phát gói trong vài microsecond, tuyệt đối không để tràn bộ đệm làm rơi gói tin!
  \item \textbf{ECN (Explicit Congestion Notification) \& DCQCN}: Switch chủ động đánh dấu 2 bit cờ ECN trong tiêu đề gói tin IP khi hàng đợi bắt đầu tích lũy. Card mạng bên nhận phát hiện cờ ECN sẽ gửi bản tin \code{CNP (Congestion Notification Packet)} yêu cầu bên phát tự động giảm nhẹ tốc độ truyền trước khi hiện tượng nghẽn mạng xảy ra, ngăn chặn hiện tượng bão hòa khung PAUSE (\textit{PFC Deadlock}).
\end{itemize}

\subsection{Giải Tích Thuật Toán All-Reduce Trong NCCL: Ring vs Tree Topology}
Thư viện truyền thông tập thể của NVIDIA (\term{NCCL - NVIDIA Collective Communications Library}) là trái tim kết nối các GPU trong quá trình suy luận và huấn luyện phân tán. Trong số các phép toán truyền thông (All-Gather, Reduce-Scatter, Broadcast), phép toán \term{All-Reduce} là phép toán quan trọng và tốn nhiều băng thông nhất: $P$ GPU, mỗi GPU nắm giữ một tensor có kích thước $S$ byte, cần cộng dồn các tensor này lại và phân phát kết quả tổng đến toàn bộ $P$ GPU.

\paragraph{1. Thuật toán Ring-AllReduce}:
Trong thuật toán Ring-AllReduce, $P$ GPU được sắp xếp thành một vòng tròn logic kín: $\text{GPU}_0 \to \text{GPU}_1 \to \dots \to \text{GPU}_{P-1} \to \text{GPU}_0$.
Tensor kích thước $S$ được chia đều thành $P$ phần bằng nhau, mỗi phần có kích thước $\frac{S}{P}$ byte.

Thuật toán gồm hai pha nối tiếp:
\begin{itemize}
  \item \textbf{Pha 1: Scatter-Reduce ($P - 1$ bước)}:
  Ở mỗi bước, mỗi GPU gửi một khối dữ liệu $\frac{S}{P}$ cho nút kế tiếp và nhận một khối từ nút liền trước, sau đó cộng dồn vào tensor cục bộ. Sau đúng $P - 1$ bước, mỗi GPU nắm giữ kết quả tổng cuối cùng của đúng một khối dữ liệu $\frac{S}{P}$.
  Lượng dữ liệu gửi đi của mỗi GPU trong pha này là:
  \begin{equation}
  \text{Bytes}_{\text{scatter}} = (P - 1) \times \frac{S}{P}
  \end{equation}

  \item \textbf{Pha 2: All-Gather ($P - 1$ bước)}:
  Mỗi GPU gửi khối dữ liệu đã được tính tổng cho nút kế tiếp để phân phát ra toàn bộ vòng. Sau đúng $P - 1$ bước nữa, toàn bộ $P$ GPU đều sở hữu bản sao hoàn chỉnh của tensor tổng.
  Lượng dữ liệu gửi đi của mỗi GPU trong pha này là:
  \begin{equation}
  \text{Bytes}_{\text{gather}} = (P - 1) \times \frac{S}{P}
  \end{equation}
\end{itemize}

Tổng lượng dữ liệu truyền qua card mạng của mỗi GPU trong toàn bộ quá trình:
\begin{equation}
S_{\text{ring}} = \text{Bytes}_{\text{scatter}} + \text{Bytes}_{\text{gather}} = 2 \times \left(\frac{P - 1}{P}\right) S
\end{equation}

\begin{theorem}[Tính Độc Lập Băng Thông Của Ring-AllReduce]
Khi số lượng GPU $P$ tăng lên rất lớn ($P \to \infty$):
\begin{equation}
\lim_{P \to \infty} 2 \left(\frac{P - 1}{P}\right) S = 2S
\end{equation}
Tổng lượng dữ liệu mà mỗi card mạng GPU phải truyền chỉ xấp xỉ gấp đôi kích thước tensor $2S$, \textbf{hoàn toàn không tăng theo số lượng GPU trong cụm}!
\end{theorem}

Thời gian thực thi của Ring-AllReduce trên mạng có độ trễ $\alpha$ và băng thông $\beta$:
\begin{equation}
T_{\text{ring}} = 2(P - 1)\alpha + 2 \left(\frac{P - 1}{P}\right) \frac{S}{\beta}
\end{equation}

\paragraph{2. Thuật toán Tree-AllReduce (Cây Nhị Phân Kép - Double Binary Tree)}:
Từ công thức trên, ta nhận thấy một điểm yếu chí mạng của Ring-AllReduce: Số bước truyền dữ liệu tỷ lệ thuận tuyến tính với số lượng GPU: $2(P - 1)$.
Khi cụm máy chủ mở rộng lên quy mô lớn (ví dụ $P = 1.024$ GPU), số bước nhảy là $2.046$ bước!
Đối với các tensor kích thước nhỏ hoặc trong pha Decode của Tensor Parallelism (nơi kích thước tensor chỉ vài chục Kilobytes), thời gian truyền bị chi phối hoàn toàn bởi độ trễ $\alpha$ (\textit{Latency-Bound}) chứ không phải băng thông $\beta$.

Để khắc phục, NCCL áp dụng thuật toán \term{Tree-AllReduce}: Dữ liệu được thu gom theo cấu trúc cây nhị phân (Reduce) lên nút gốc, sau đó phát tán ngược lại (Broadcast) xuống các nhánh lá.
Thời gian thực thi của Tree-AllReduce:
\begin{equation}
T_{\text{tree}} = 2 \log_2(P) \alpha + 2 \frac{S}{\beta_{\text{tree}}}
\end{equation}
Số bước nhảy giảm ngoạn mục từ $2(P - 1)$ xuống chỉ còn $2 \log_2(P)$! Với $P = 1.024$ GPU, số bước nhảy giảm từ $2.046$ bước xuống chỉ còn $2 \times 10 = 20$ bước, giúp tăng tốc độ trao đổi tensor nhỏ lên hàng chục lần.

---

\section{Hệ Thống Bài Tập \& Câu Hỏi Tư Duy Đa Tầng}
\label{sec:ch10-exercises}

\subsection{Nhóm A: Câu Hỏi Tư Duy Bản Chất (Conceptual Reasoning)}
\begin{enumerate}
  \item \textbf{Câu A.1}: Tại sao việc chia sẻ dung lượng bộ nhớ IPC (\code{--shm-size}) lại cần thiết cho cụm GPU chạy Tensor Parallelism sử dụng Ray, trong khi một ứng dụng Web API thông thường chỉ cần vài chục megabytes RAM?
  \item \textbf{Câu A.2}: Phân tích nguyên nhân tại sao Pipeline Parallelism (PP) tạo ra hiện tượng ``Bong bóng tính toán'' (Execution Bubbles), và tại sao tăng số lượng Micro-batches lại giúp giảm tỷ lệ phần trăm thời gian chết này.
  \item \textbf{Câu A.3}: Tại sao trong một image Docker phục vụ sản xuất, chúng ta tuyệt đối không nên cài đặt bộ công cụ đầy đủ \code{cuda-toolkit} (chứa trình biên dịch \code{nvcc}), mà chỉ nên cài đặt gói \code{cuda-runtime}?
  \item \textbf{Câu A.4}: Giải thích tại sao cơ chế Horizontal Pod Autoscaler (HPA) chuẩn của Kubernetes dựa trên mức sử dụng CPU/Memory lại hoàn toàn thất bại khi áp dụng cho các cụm phục vụ vLLM.
  \item \textbf{Câu A.5}: Phân biệt sự khác nhau giữa Continuous Batching trong vLLM và Dynamic Batching trong Triton Inference Server. Khi nào một hệ thống nên ưu tiên chọn Triton thay vì vLLM?
  \item \textbf{Câu A.6}: Tại sao mạng InfiniBand và RoCEv2 lại bắt buộc cơ chế điều khiển luồng không mất gói (Lossless Flow Control), trong khi các dịch vụ Web API Internet vẫn hoạt động hoàn hảo trên mạng TCP có tỷ lệ mất gói $0.1\%-1\%$?
  \item \textbf{Câu A.7}: Trong thuật toán Ring-AllReduce, tại sao tổng lượng dữ liệu truyền qua mỗi card mạng lại độc lập với số lượng GPU $P$ khi $P \to \infty$? Khi nào thì thuật toán Tree-AllReduce vượt trội hơn Ring-AllReduce?
\end{enumerate}

\subsection{Nhóm B: Bài Toán Tính Toán Giải Tích Bằng Tay (Analytical Problems)}
\begin{enumerate}
  \item \textbf{Bài B.1 (Tính Định Cỡ Băng Thông Phục Vụ Quy Mô 1000 Người Dùng)}: \\
  Hệ thống suy luận phục vụ mô hình Llama-3-8B ($d_{\text{model}} = 4.096$, $n_{\text{layers}} = 32$) sử dụng Tensor Parallelism với $TP = 2$ GPU A100 PCIe.
  Giả sử hệ thống phục vụ tải cực đại $B = 256$ người dùng đồng thời trong pha Decode.
  \begin{itemize}
    \item[a)] Tính lượng dữ liệu All-Reduce cần truyền cho mỗi token được sinh ra.
    \item[b)] Nếu băng thông PCIe thực tế giữa 2 card là $32\text{ GB/s}$, tính thời gian truyền thông mạng tối thiểu cho mỗi token và đánh giá xem hệ thống có thể duy trì tốc độ $50\text{ tokens/giây}$ cho người dùng hay không.
  \end{itemize}

  \item \textbf{Bài B.2 (Hiệu Suất Pipeline Bubbles Trong 1F1B Scheduling)}: \\
  Xét một hệ thống Pipeline Parallelism gồm $K = 4$ tầng (4 stages).
  Hệ thống chia lô dữ liệu thành $M$ micro-batches.
  Tỷ lệ thời gian lãng phí do bong bóng Pipeline được tính theo công thức:
  \begin{equation*}
  F_{\text{bubble}} = \frac{K - 1}{M + K - 1}
  \end{equation*}
  Hãy tính tỷ lệ lãng phí $F_{\text{bubble}}$ khi $M = 4$ micro-batches, và khi $M = 32$ micro-batches. Rút ra bài học kỹ thuật khi chọn kích thước batch.

  \item \textbf{Bài B.3 (Định Cỡ KEDA Pod Replica Cho Sự Kiện Lưu Lượng Đỉnh)}: \\
  Vào ngày Black Friday, một sàn thương mại điện tử ghi nhận $1.200$ yêu cầu phân tích đánh giá sản phẩm đồng thời.
  Mỗi GPU Pod phục vụ mô hình phân tích có dung lượng xử lý tối ưu là $C_{\text{pod}} = 40$ requests đồng thời.
  Cấu hình KEDA có giới hạn $R_{\min} = 2$ và $R_{\max} = 50$.
  Tính số lượng Pod mà KEDA sẽ kích hoạt. Nếu chi phí thuê một máy chủ GPU là $2.5\text{ USD/giờ}$, tính chi phí phần cứng trong 4 giờ cao điểm.

  \item \textbf{Bài B.4 (So Sánh Định Lượng Ring-AllReduce vs Tree-AllReduce Trên Cụm 64 GPU)}: \\
  Xét một cụm gồm $P = 64$ máy chủ GPU H100 kết nối qua mạng InfiniBand $400\text{ Gbps}$ ($\beta = 50\text{ GB/s}$), độ trễ mỗi chặng mạng $\alpha = 1.5\mu\text{s}$.
  So sánh thời gian truyền thông giữa Ring-AllReduce và Tree-AllReduce trong hai trường hợp:
  \begin{itemize}
    \item[a)] Tensor nhỏ: $S = 64\text{ KB}$ (trong pha Decode của Tensor Parallelism).
    \item[b)] Tensor lớn: $S = 1\text{ GB}$ (đồng bộ gradient trong Data Parallelism). Rút ra kết luận về chiến lược chuyển mạch cấu trúc của thư viện NCCL.
  \end{itemize}

  \item \textbf{Bài B.5 (Tính Định Cỡ Bộ Đệm Switch Headroom Ngăn Chặn Mất Gói PFC)}: \\
  Một cổng Switch Ethernet $400\text{ Gbps}$ nối tới card mạng H100 qua cáp quang dài $L = 50\text{ mét}$ (tốc độ truyền tín hiệu trong sợi quang $v = 2 \times 10^8\text{ m/s}$).
  Thời gian phát hiện nghẽn và phản hồi khung PAUSE của phần cứng Switch và NIC là $t_{\text{proc}} = 400\text{ ns}$.
  Tính dung lượng bộ đệm tối thiểu (Headroom Buffer) cần dự trữ trên cổng switch để đảm bảo không một bit dữ liệu nào bị tràn rơi khi kích hoạt cơ chế PFC.
\end{enumerate}

\subsection{Nhóm C: Thực Hành Kỹ Nghệ \& Lập Trình (Implementation Labs)}
\begin{enumerate}
  \item \textbf{Bài C.1 (Xây Dựng Cụm Docker Compose Giám Sát vLLM)}: \\
  Viết tệp cấu hình \code{docker-compose.yml} hoàn chỉnh thiết lập cụm dịch vụ:
  \begin{itemize}
    \item Service 1: Khởi chạy \code{vllm-server} nạp mô hình, cấu hình giới hạn GPU qua trường tài nguyên Compose, mở cổng metrics 8000.
    \item Service 2: Khởi chạy \code{prometheus} cào dữ liệu định kỳ 5 giây một lần từ đường dẫn giám sát của vLLM.
    \item Service 3: Khởi chạy \code{grafana} tự động nạp Dashboard theo dõi TTFT và tỷ lệ chiếm dụng bộ nhớ KV Cache.
  \end{itemize}

  \item \textbf{Bài C.2 (Hiện Thực Hóa Bộ Lập Lịch Dynamic Batch Scheduler An Toàn)}: \\
  Cài đặt hoàn chỉnh lớp Python \code{DynamicBatchScheduler} với đầy đủ cơ chế kiểm tra timeout và kích thước lô cực đại, kèm theo unit test kiểm tra tính chính xác của thứ tự xuất xưởng các request.

  \item \textbf{Bài C.3 (Tự Động Sinh Tệp Cấu Hình Triton config.pbtxt)}: \\
  Viết một hàm Python tự động nhận vào các thông số kích thước tensor đầu vào, kiểu dữ liệu, và kích thước batch mong muốn, xuất ra nội dung tệp \code{config.pbtxt} chuẩn cú pháp Protobuf của NVIDIA Triton.
\end{enumerate}

\subsection{Nhóm D: Phân Tích Thất Bại \& Debugging Thực Tế (Post-Mortem Analysis)}
\begin{enumerate}
  \item \textbf{Tình huống sự cố 1 (Treo cứng mạng NCCL trong Tensor Parallelism)}: \\
  Kỹ sư triển khai mô hình Llama-3-70B trên máy chủ 8 GPU H100 với tham số \code{--tensor-parallel-size 8}.
  Khi tiến trình khởi động, log hiển thị việc nạp trọng số thành công $100\%$. Tuy nhiên, ngay tại bước khởi tạo nhóm giao tiếp NCCL (\code{ncclCommInitRank}), toàn bộ 8 tiến trình GPU bị treo cứng vĩnh viễn (Hang/Freeze), mức sử dụng CPU rơi về $0\%$, và không có bất kỳ thông báo lỗi nào xuất hiện.
  Hãy phân tích nguyên nhân gốc rễ (Root Cause) và chỉ ra cách debug giải quyết lỗi treo mạng NCCL này.

  \item \textbf{Tình huống sự cố 2 (Tràn ổ đĩa đệm OverlayFS do ghi đệm mô hình)}: \\
  Sau 2 ngày chạy dịch vụ vLLM trên cụm Kubernetes, một loạt máy chủ vật lý đồng loạt báo động đỏ: ổ đĩa gốc \code{/} bị chiếm dụng $100\%$ dung lượng ($500\text{ GB}$). Khi kiểm tra, thư mục lưu trữ Docker Overlay2 phình to bất thường vượt tầm kiểm soát. Phân tích nguyên nhân và thiết lập cấu hình chuẩn để ngăn chặn sự cố này tái diễn.
\end{enumerate}

---

\section{Hướng Dẫn Lời Giải \& Bình Luận Kỹ Thuật Chi Tiết}
\label{sec:ch10-solutions}

\subsection{Lời Giải Nhóm A: Câu Hỏi Tư Duy Bản Chất}
\begin{itemize}
  \item \textbf{Lời giải Câu A.1}:
  Trong kiến trúc Tensor Parallelism, mô hình được chia nhỏ thành nhiều tiến trình độc lập (mỗi tiến trình quản lý 1 GPU). Để chuyển giao các tensor dữ liệu lớn giữa các tiến trình CPU (Host Processes) mà không phải sao chép qua giao thức mạng TCP loopback chậm chạp, PyTorch và Ray sử dụng cơ chế IPC Shared Memory (Inter-Process Communication qua \code{/dev/shm}). Cơ chế này ánh xạ cùng một dải bộ nhớ RAM vật lý vào không gian địa chỉ của cả 8 tiến trình. Một tensor kích hoạt cho batch lớn có thể chiếm vài trăm megabytes; nếu chỉ dùng $64\text{ MB}$ mặc định của Docker, phân vùng này bị cạn kiệt ngay lập tức, dẫn tới lỗi \code{SIGBUS}.

  \item \textbf{Lời giải Câu A.2}:
  Trong Pipeline Parallelism, Stage 2 không thể tính toán cho đến khi Stage 1 hoàn tất việc xử lý dữ liệu và gửi Activation qua. Do đó, trong giai đoạn khởi động đường ống (Warm-up phase), Stage 2, Stage 3 và Stage 4 hoàn toàn phải ngồi rảnh rỗi. Khi ta chia một batch lớn thành $M$ micro-batches rất nhỏ, Stage 1 xử lý xong micro-batch 1 là lập tức gửi ngay cho Stage 2, đồng thời Stage 1 bắt tay vào xử lý micro-batch 2. Bằng cách tăng $M$, thời gian các stage cùng làm việc đồng thời chiếm đại đa số thời gian tổng thể, làm cho tỷ lệ thời gian chết $\frac{K-1}{M + K - 1}$ tiệm cận về $0$.

  \item \textbf{Lời giải Câu A.3}:
  Gói \code{cuda-toolkit} chứa đầy đủ trình biên dịch mã C++ \code{nvcc}, các tệp header (\code{.h}), và các công cụ phát triển chuyên sâu có dung lượng lên tới $4\text{ GB} - 6\text{ GB}$. Trong môi trường sản xuất (Production Runtime), mô hình chỉ cần các thư viện liên kết động (\code{.so}) để thực thi mã máy đã được biên dịch sẵn. Cài đặt toolkit không chỉ làm kích thước Docker Image phình to gấp đôi (làm chậm thời gian kéo ảnh trên cụm Kubernetes), mà còn vi phạm nguyên lý an ninh mạng (tăng bề mặt tấn công khi để sẵn trình biên dịch mã trên máy chủ sản xuất).

  \item \textbf{Lời giải Câu A.4}:
  vLLM chiếm dụng trước (Pre-allocate) gần như toàn bộ VRAM rảnh rỗi cho bộ đệm PagedAttention ngay khi nạp mô hình để tối ưu hóa bộ nhớ, do đó chỉ số bộ nhớ luôn đạt $\sim 90\%-95\%$ bất kể hệ thống đang rảnh rỗi hay có tải nặng. Đồng thời, CPU chỉ thực hiện nhiệm vụ điều phối luồng bất đồng bộ nhẹ nên mức sử dụng CPU luôn ở mức thấp. HPA chuẩn chỉ nhìn vào hai chỉ số này sẽ hoàn toàn bị mù trước tình trạng ứ đọng hàng trăm request đang xếp hàng chờ.

  \item \textbf{Lời giải Câu A.5}:
  Continuous Batching của vLLM hoạt động ở cấp độ bước sinh token (Iteration-level): mỗi khi một request sinh xong token kết thúc, nó lập tức được đưa ra khỏi lô và một request mới từ hàng đợi được chèn vào ngay tại bước tiếp theo mà không cần chờ toàn bộ lô hoàn thành. Dynamic Batching của Triton gom các request độc lập theo cửa sổ thời gian timeout. Triton được ưu tiên khi hệ thống cần xây dựng pipeline đa mô hình phức tạp (ví dụ: Speech-to-Text $\to$ LLM $\to$ Rerank) hoặc khi cần phục vụ song song các mô hình học máy truyền thống và thị giác máy tính.

  \item \textbf{Lời giải Câu A.6}:
  Trong mạng TCP thông thường, cơ chế SACK hoặc Fast Retransmit cho phép gửi lại các gói tin thất lạc; độ trễ vài chục mili-giây không ảnh hưởng nghiêm trọng đến người dùng web. Tuy nhiên trong cụm AI sử dụng RDMA (RoCEv2), card mạng không có bộ đệm lớn và phần cứng áp dụng thuật toán Go-Back-N. Khi rơi một gói tin duy nhất, toàn bộ khối truyền hàng trăm Megabyte bị hủy bỏ và phải truyền lại từ đầu. Hơn nữa, vì các GPU phải đồng bộ kết quả tại mỗi bước tính toán qua rào chắn (Barrier), việc một GPU bị chậm do mất gói tin sẽ khiến toàn bộ hàng ngàn GPU trong cụm phải dừng lại chờ đợi, làm sụp đổ thông lượng hệ thống. Do đó, mạng AI bắt buộc phải là Lossless Ethernet với cơ chế PFC và ECN.

  \item \textbf{Lời giải Câu A.7}:
  Trong thuật toán Ring-AllReduce, lượng dữ liệu mỗi GPU truyền đi qua hai pha là $2 \frac{P-1}{P} S$. Khi số lượng GPU $P \to \infty$, tỷ số $\frac{P-1}{P} \to 1$, do đó lượng dữ liệu truyền tiệm cận hằng số $2S$, hoàn toàn không phụ thuộc vào kích thước cụm.
  Tuy nhiên, số bước truyền là $2(P-1)$. Khi kích thước tensor $S$ nhỏ (ví dụ vài chục Kilobytes trong pha Decode), thời gian truyền bị chi phối bởi độ trễ $\alpha$ qua $2(P-1)$ bước nhảy. Tree-AllReduce tổ chức phân tán theo cây nhị phân kép chỉ mất $2\log_2(P)$ bước, giúp giảm độ trễ trao đổi tensor nhỏ lên hàng chục lần so với Ring-AllReduce.
\end{itemize}

\subsection{Lời Giải Nhóm B: Bài Toán Tính Toán Giải Tích Bằng Tay}

\paragraph{Lời giải Bài B.1}:
\begin{itemize}
  \item[a)] Kích thước tensor cho mỗi phép All-Reduce:
  \begin{equation*}
  S = B \times d_{\text{model}} \times b = 256 \times 4.096 \times 2\text{ Bytes} = 2.097.152\text{ Bytes} = 2\text{ MiB}
  \end{equation*}
  Với $TP = 2$, mỗi phép All-Reduce truyền:
  \begin{equation*}
  D_{\text{op}} = 2 \times \left( \frac{2 - 1}{2} \right) \times 2\text{ MiB} = 2\text{ MiB}
  \end{equation*}
  Mỗi tầng Transformer có 2 phép All-Reduce, với 32 tầng:
  \begin{equation*}
  D_{\text{total}} = 32 \times 2 \times 2\text{ MiB} = 128\text{ MiB} \approx 0.1342\text{ GB}
  \end{equation*}

  \item[b)] Thời gian truyền thông tối thiểu qua bus PCIe ($32\text{ GB/s}$):
  \begin{equation*}
  t_{\text{comm}} = \frac{0.1342\text{ GB}}{32\text{ GB/s}} \approx 0.00419\text{ giây} \approx 4.19\text{ ms}
  \end{equation*}
  Để đạt tốc độ $50\text{ tokens/giây}$ cho người dùng, tổng thời gian cho mỗi token (bao gồm cả tính toán và truyền thông) phải nhỏ hơn:
  \begin{equation*}
  T_{\text{max}} = \frac{1}{50} = 20\text{ ms}
  \end{equation*}
  Vì $t_{\text{comm}} \approx 4.19\text{ ms} < 20\text{ ms}$, hệ thống hoàn toàn có thể đáp ứng được nếu thời gian tính toán thực thi nhân ma trận trên GPU nhỏ hơn $20 - 4.19 = 15.81\text{ ms}$. Tuy nhiên, thời gian giao tiếp đã chiếm tới hơn $20\%$ ngân sách độ trễ!
\end{itemize}

\paragraph{Lời giải Bài B.2}:
Áp dụng công thức tỷ lệ lãng phí do bong bóng Pipeline $F_{\text{bubble}} = \frac{K - 1}{M + K - 1}$ với $K = 4$:
\begin{itemize}
  \item Khi $M = 4$ micro-batches:
  \begin{equation*}
  F_{\text{bubble}} = \frac{4 - 1}{4 + 4 - 1} = \frac{3}{7} \approx 42.86\%
  \end{equation*}
  Gần một nửa thời gian phần cứng GPU bị bỏ phí ngồi chờ!

  \item Khi $M = 32$ micro-batches:
  \begin{equation*}
  F_{\text{bubble}} = \frac{4 - 1}{32 + 4 - 1} = \frac{3}{35} \approx 8.57\%
  \end{equation*}
  Tỷ lệ lãng phí giảm xuống dưới $9\%$, hiệu suất sử dụng GPU tăng vọt lên hơn $91\%$.
\end{itemize}

\paragraph{Lời giải Bài B.3}:
\begin{itemize}
  \item Số lượng Pod mong muốn:
  \begin{equation*}
  R = \left\lceil \frac{1.200}{40} \right\rceil = 30\text{ Pods}
  \end{equation*}
  Vì $R_{\min} = 2 \le 30 \le R_{\max} = 50$, KEDA sẽ kích hoạt đúng $30$ Pods.
  \item Chi phí phần cứng trong 4 giờ:
  \begin{equation*}
  \text{Cost} = 30\text{ Pods} \times 2.5\text{ USD/Pod/giờ} \times 4\text{ giờ} = 300\text{ USD}
  \end{equation*}
\end{itemize}

\paragraph{Lời giải Bài B.4}:
\begin{itemize}
  \item[a)] Với tensor nhỏ $S = 64\text{ KB} = 6.4 \times 10^{-5}\text{ GB}$:
  \begin{itemize}
    \item Ring-AllReduce: Số bước truyền là $2(64 - 1) = 126$ bước.
    Thời gian độ trễ: $126 \times 1.5\mu\text{s} = 189\mu\text{s}$.
    Thời gian băng thông: $2 \times \frac{63}{64} \times \frac{6.4 \times 10^{-5}\text{ GB}}{50\text{ GB/s}} \approx 2.52\mu\text{s}$.
    Tổng thời gian Ring: $T_{\text{ring}} = 189 + 2.52 \approx \mathbf{191.52\mu\text{s}}$.
    \item Tree-AllReduce: Số bước truyền là $2 \log_2(64) = 2 \times 6 = 12$ bước.
    Thời gian độ trễ: $12 \times 1.5\mu\text{s} = 18\mu\text{s}$.
    Thời gian băng thông (băng thông hiệu dụng cây $\beta/2 = 25\text{ GB/s}$): $2 \times \frac{6.4 \times 10^{-5}}{25} \approx 5.12\mu\text{s}$.
    Tổng thời gian Tree: $T_{\text{tree}} = 18 + 5.12 \approx \mathbf{23.12\mu\text{s}}$.
    $\implies$ \textbf{Tree-AllReduce nhanh hơn hơn 8 lần!}
  \end{itemize}

  \item[b)] Với tensor lớn $S = 1\text{ GB}$:
  \begin{itemize}
    \item Ring-AllReduce: Thời gian băng thông: $2 \times \frac{63}{64} \times \frac{1\text{ GB}}{50\text{ GB/s}} \approx 0.03938\text{s} = 39.38\text{ms}$.
    Tổng thời gian Ring: $T_{\text{ring}} = 0.189\text{ms} + 39.38\text{ms} \approx \mathbf{39.57\text{ms}}$.
    \item Tree-AllReduce: Thời gian băng thông: $2 \times \frac{1\text{ GB}}{25\text{ GB/s}} = 0.08\text{s} = 80\text{ms}$.
    Tổng thời gian Tree: $T_{\text{tree}} = 0.018\text{ms} + 80\text{ms} \approx \mathbf{80.02\text{ms}}$.
    $\implies$ \textbf{Ring-AllReduce nhanh gấp đôi Tree-AllReduce!}
  \end{itemize}
  \textbf{Kết luận}: NCCL tự động áp dụng Tree-AllReduce cho tensor nhỏ (pha Decode) và chuyển sang Ring-AllReduce cho tensor lớn (pha Prefill và All-Reduce gradient).
\end{itemize}

\paragraph{Lời giải Bài B.5}:
Tổng thời gian kể từ khi switch gửi khung PAUSE đến khi dòng gói tin trên cổng switch dừng hẳn:
\begin{enumerate}
  \item Thời gian khung PAUSE truyền trên cáp quang tới card mạng:
  \begin{equation*}
  t_{\text{prop}} = \frac{L}{v} = \frac{50\text{ m}}{2 \times 10^8\text{ m/s}} = 2.5 \times 10^{-7}\text{s} = 250\text{ ns}
  \end{equation*}
  \item Thời gian card mạng phân tích khung PAUSE và dừng bộ điều khiển MAC: $t_{\text{proc}} = 400\text{ ns}$.
  \item Các gói tin đã rời card mạng trước khi dừng tiếp tục bay tới switch: $t_{\text{prop}} = 250\text{ ns}$.
\end{enumerate}
Tổng thời gian xả dữ liệu tiếp diễn:
\begin{equation*}
T_{\text{drain}} = 2 \times t_{\text{prop}} + t_{\text{proc}} = 2(250\text{ ns}) + 400\text{ ns} = 900\text{ ns} = 0.9\mu\text{s}
\end{equation*}
Băng thông đường truyền là $400\text{ Gbps} = 50\text{ GB/s} = 50 \times 10^9\text{ Bytes/giây}$.
Dung lượng bộ đệm tối thiểu cần thiết để không rơi gói tin:
\begin{equation*}
\text{Headroom} = T_{\text{drain}} \times \text{Bandwidth} = 900 \times 10^{-9}\text{s} \times 50 \times 10^9\text{ Bytes/s} = \mathbf{45.000\text{ Bytes}} \approx \mathbf{43.95\text{ KiB}}
\end{equation*}

\subsection{Lời Giải Nhóm C: Thực Hành Kỹ Nghệ \& Lập Trình}
Cấu hình `docker-compose.yml` chuẩn mực và mã nguồn Python đã được trình bày chi tiết trong Mục~\ref{sec:production-docker-keda} và Mục~\ref{sec:python-serving-modules}.

\subsection{Lời Giải Nhóm D: Phân Tích Thất Bại \& Debugging Thực Tế}

\paragraph{Lời giải Sự cố 1 (Treo cứng mạng NCCL)}:
\begin{itemize}
  \item \textbf{Nguyên nhân gốc rễ}: NCCL mặc định quét toàn bộ card mạng và nhầm lẫn chọn card mạng ảo Docker (\code{docker0}) thay vì card mạng vật lý hoặc card InfiniBand, khiến các GPU không thể bắt tay với nhau.
  \item \textbf{Khắc phục}: Đặt biến môi trường \code{export NCCL_SOCKET_IFNAME=eth0,ib0} và chạy container với \code{--ipc=host --net=host}.
\end{itemize}

\paragraph{Lời giải Sự cố 2 (Tràn ổ đĩa OverlayFS)}:
\begin{itemize}
  \item \textbf{Nguyên nhân gốc rễ}: Tiến trình tải mô hình từ Hugging Face Hub trực tiếp vào thư mục gốc mặc định của container (\code{~/.cache/huggingface/}). Do cơ chế Copy-on-Write của OverlayFS, toàn bộ trọng số được ghi vào tầng tạm \code{UpperDir} của container trên ổ đĩa hệ điều hành, làm cạn kiệt đĩa.
  \item \textbf{Khắc phục}: Luôn gắn kết một Persistent Volume chuyên dụng vào thư mục cache:
  \begin{lstlisting}[language=bash]
docker run -v /mnt/fast_nvme/hf_cache:/app/cache/huggingface \
           -e HF_HOME=/app/cache/huggingface ...
  \end{lstlisting}
\end{itemize}
